\subsection{The proper holomorphic family and coherent descent}\hypertarget{tag-01PC}{}\label{fam:section:A.3}

\begin{proposition}[Coherent even shifts and the quotient family]\hypertarget{tag-01Q2}{}\label{fam:src:A.3.1}\label{fam:st:A.3.1}
\StatementRecord{fam:src:A.3.1}{proved}{01Q2}{\SeriesRef{fam:src:A.2.1}}

For each $n\in2\mathbb Z$ there is a biholomorphism $F_{n,c}:X_c\to X_{c+n}$, holomorphic jointly in $c$, with $F_{2h,c+2k}F_{2k,c}=F_{2(h+k),c}$. The quotient family over $0<|w|<e^{-\pi M}$, $w=e^{-\pi ic}$, is proper and holomorphic, with pairwise non-biholomorphic fibres.

Before the quotient, every disc $Q$ with $\overline Q\subset\mathcal P$
has the proper holomorphic submersion and proper base map
\[
 \pi:\mathcal X_Q\longrightarrow Q,\qquad
 F:\mathcal X_Q\longrightarrow B\times Q,\qquad
 \pi=\operatorname{pr}_Q\circ F.
 \tag{PD2}
\]
Their restrictions to $c$ are $X_c$ and $f_c$, and their three
distinguished values are the constant sections $0,1,\infty$.
\end{proposition}
\begin{proof}[Proof of the proper-family assertion; coherent descent in \cref{I:coherent-descent-proof}]
\begin{enumerate}[label=\arabic*.,leftmargin=*]
\item \emph{The holomorphic family.}
Let $Q$ be a disc with $\overline Q\subset\mathcal P$.
On an ordinary lifted base chart use the period matrix
\eqref{wb:pshift:eq:periods}, jointly holomorphic in $(z,c)$.
Its real determinant is $\Im\tau\,D_c$, and its inverse is
uniformly bounded on compact base-parameter sets. Thus the lattice
action is free and properly discontinuous there. A fixed closed
real lattice parallelotope gives compact surjective representatives
over each compact subset of the base chart times $Q$.
The quotient is holomorphic and its projection to $Q$ is a submersion.
The covariance convention of \cref{I:def:affine-covariance} is
\[
 R_g(z)\Pi(z,c)=\Pi(gz,c)A_g,
 \qquad M_g=A_g^{-1},\qquad
 \Pi(gz,c)=R_g(z)\Pi(z,c)M_g.
\]
Consequently the base deck quotient is jointly holomorphic; its
proper discontinuity follows from that of the lifted base action.

At a finite root chart the family action is
\[
 (z,c,\zeta)\longmapsto
 \left(g_jz,c,R_{g_j}(z)\zeta+
             \Pi(g_jz,c)\frac{v_j}{m_j}\right),
 \qquad (m_1,m_2)=(3,4).
 \tag{PD7}
\]
In the real marking it is $(z,c,u)\mapsto(g_jz,c,A_ju+v_j/m_j)$.
The invariant-covector freeness test in \cref{fam:src:A.1.2}
is independent of $c$, so these quotients are local biholomorphisms
and submersions to $Q$. The same parallelotope proves properness
over the finite base disc times $Q$. Their punctured connectors are
\[
 \frac{\log s_j(z)}{2\pi i}\Pi(z,c)v_j.
 \tag{PD8}
\]
A logarithm change adds a period, so this is a jointly holomorphic
torus section. Writing $\sigma_j=(\log s_j)/(2\pi i)\,\Pi(z,c)v_j$, the
relations $A_jv_j=v_j$ and
$\log s_j(g_jz)=\log s_j(z)-2\pi i/m_j$ give
\[
 R_{g_j}(z)\sigma_j(z,c)-\sigma_j(g_jz,c)
       =\Pi(g_jz,c)v_j/m_j.
\]
Thus translation by $\sigma_j$ intertwines the finite affine action
with the middle linear action, with the specified root coordinate.

\item \emph{Uniform cusp charts and properness.}
On the fixed unimodular fan write the full cusp period matrix as
\[
 C(t_c,c)=\begin{pmatrix}6\mu(t_c)&h(t_c)\\
 b^{\rm part}(t_c)-h(t_c)+c&\mu(t_c)\end{pmatrix},
 \qquad B_0=\begin{pmatrix}0&1\\-1&0\end{pmatrix}.
 \tag{PD9}
\]
Choose a closed cusp disc $|t_c|\le r_1$ on which these functions
are holomorphic, and set
\[
 K_Q=\max_{|t_c|\le r_1,\ c\in\overline Q}
               \|-2\pi\Im C(t_c,c)\|,
 \qquad 0<\epsilon<\min(r_1,1),\quad
               |\log\epsilon|>2K_Q.
 \tag{PD10}
\]
The lattice acts on the toric family times $Q$ by
\[
 (x,t_c,c)\longmapsto
 \bigl(e^{2\pi iC(t_c,c)\lambda}t_c^{B_0\lambda}x,t_c,c\bigr),
 \qquad\lambda\in\mathbb Z^2.
 \tag{PD11}
\]
The fan translation and nowhere-zero torus multiplier extend this
map across every toric stratum. With Euclidean operator norm,
$B_0^{\mathsf t}B_0=I$, and
\[
 B_{t_c,c}=B_0+\frac{-2\pi\Im C(t_c,c)}{\log|t_c|},
 \quad \|B_{t_c,c}-B_0\|<\tfrac12,
 \quad \tfrac12\|\lambda\|<\|B_{t_c,c}\lambda\|
                    \le\tfrac32\|\lambda\|\quad(\lambda\ne0).
 \tag{PD12}
\]
These inequalities exclude stabilizers on the dense torus.
For a height-one triangle, use its three dual unimodular cone
coordinates $z_0,z_1,z_2$, with $t_c=z_0z_1z_2$.
In the open chart $|z_i|<2$, $|t_c|<\epsilon$, the normalized
logarithmic barycentric coefficients have sum one and are bounded
below by $-\log2/|\log\epsilon|$. Thus the logarithmic positions
in each fixed chart lie in a bounded set, uniformly in $c$.
For a fixed pair of charts, (PD12) bounds the finitely many
$\lambda$ which can translate one into the other. At the central
fibre a nonzero $B_0\lambda$ moves every bounded fan cell.
These charts cover the entire toric space, including its boundary.
On the dense torus the duality identity is
\[
 \log|z_i|=\log|t_c|\,\theta_i^\sigma(y),
\]
where $\theta_i^\sigma$ are the barycentric coordinates of the
height-one triangle. A triangle containing $y$ therefore gives
$|z_i|\le1$. For a boundary point in the orbit of the cone over a
cell $P$, the coordinates belonging to vertices of $P$ vanish.
The remaining characters are dual to the image cones of triangles
containing $P$ in the quotient lattice by $\operatorname{span}(P,1)$.
Those cones cover the quotient real space: lift a vector, add a
sufficiently small positive multiple of it to a relative interior
point of $(P,1)$, and rescale to height one. Local finiteness places
that point in a triangle containing $P$. Applied to the negative
logarithm of the orbit coordinates, this gives the same closed
unit-polydisc bound. Thus the open radius-two charts cover the
boundary as well. Every compact set has a finite subcover.
The chart-pair estimate applies to an intersection meeting the
boundary because that intersection is open and meets the dense
torus. It therefore proves freeness and proper discontinuity on
the entire toric family.

For properness, fix $\epsilon_1<\epsilon$ and let $K_0$ be the
finite union of closed unit polydiscs for height-one triangles
meeting the closed radius-$\sqrt2$ ball, intersected with
$|t_c|\le\epsilon_1$. For $t_c\ne0$, choose $\lambda$ with
$\|B_{t_c,c}^{-1}y+\lambda\|_\infty\le1/2$.
The translated logarithmic position has norm at most
$3/(2\sqrt2)<\sqrt2$. A triangle containing it has nonnegative
barycentric coordinates of sum one, so each corresponding cone
coordinate has modulus at most one. This gives a representative
in $K_0$. At $t_c=0$, translate the height-one component and use
the six closed toric bidiscs covering $E_0\simeq\mathrm dP_6$.
The multiplier $\exp(2\pi iC(0,c)B_0v)$ is included in this
translation; the bidiscs cover the resulting whole compact component.
For each compact $K\subset Q$, the compact set $K_0\times K$
therefore maps onto the full inverse image of
$\{|t_c|\le\epsilon_1\}\times K$. This proves properness.
The quotient charts are open subsets of a smooth toric variety
times $Q$, so projection to $Q$ is a submersion.

\item Exponentiating the logarithmic coordinates gives the cusp
connector with period identity $Z=sB_0+C(t_c,c)$.
Glue it and (PD8) to the middle family. The distinguished base
discs are disjoint. Points with distinct $(t,c)$ values separate
on the base; points with the same value lie in one Hausdorff local
piece. The resulting total space is therefore Hausdorff and smooth.
The local proper maps give a proper map $F$ to $B\times Q$:
cover a compact subset by finitely many relatively compact sets
inside the displayed base pieces. Compactness of $B$ proves
properness of $\pi=\operatorname{pr}_QF$.
The fibres and overlaps are the specified $X_c$ and $f_c$.
These local families agree on smaller cusp radii and glue over
$\mathcal P$.
\end{enumerate}
\end{proof}

